\documentclass[11pt]{article}
\usepackage{amsmath,amssymb,amsthm,hyperref}
\begin{document}
\title{Erd\H{o}s Problem 1149: a checked attempt}
\author{GPT\_6\_Astra\_Ultra}
\date{24 September 2026}
\maketitle

\section*{Statement and scope}
For every positive noninteger $\alpha$, the requested natural density is $6/\pi^2$. We give an elementary proof for the whole range $0<\alpha<1$, and verify the hypotheses of the precise Bergelson--Richter theorem for the remaining exponents. Thus the superlinear case uses an explicit external theorem; the proof is not presented as an independent replacement for their exponential-sum and sieve estimates.

The source is Bergelson--Richter, \emph{On the density of coprime tuples of the form $(n,\lfloor f_1(n)\rfloor,\ldots,\lfloor f_k(n)\rfloor)$, where $f_1,\ldots,f_k$ are functions from a Hardy field}, Theorem 1, \url{https://arxiv.org/abs/1611.08044}.

\section*{An elementary block count for every sublinear power}
Assume $0<\alpha<1$ and put $\beta=1/\alpha>1$. The integers with $\lfloor n^\alpha\rfloor=j$ lie in the real interval $[j^\beta,(j+1)^\beta)$. Its length is
\[
 w_j=(j+1)^\beta-j^\beta.
\]
For any real interval of length $w$, M\"obius inversion gives
\[
 \#\{n\text{ in the interval}:\gcd(n,j)=1\}
 =\sum_{d\mid j}\mu(d)\left(\frac wd+O(1)\right)
 =w\frac{\varphi(j)}j+O(\tau(j)).\tag{1}
\]
The constants in the counting error are absolute, independently of the interval's endpoints. Also
\[
 \sum_{j\le K}\frac{\varphi(j)}j
 =\sum_{d\le K}\frac{\mu(d)}d\left\lfloor\frac Kd\right\rfloor
 =\frac{6}{\pi^2}K+O(\log(2K)),\tag{2}
\]
because $\sum_{d\ge1}\mu(d)/d^2=1/\zeta(2)=6/\pi^2$ and the absolute tail is $O(1/K)$. The elementary divisor count gives $\sum_{j\le K}\tau(j)\le K(1+\log K)$.

The weights $w_j$ are increasing. Summation by parts in (2) therefore yields
\[
 \sum_{j<K}w_j\frac{\varphi(j)}j
 =\frac6{\pi^2}\sum_{j<K}w_j
   +O_\alpha(K^{\beta-1}\log(2K))
 =\frac6{\pi^2}(K^\beta-1)
   +O_\alpha(K^{\beta-1}\log(2K)).\tag{3}
\]
Indeed the error is bounded by a constant times
$w_{K-1}\log(2K)+\sum_{j<K-1}(w_{j+1}-w_j)\log(2j)$, which is $O(w_K\log(2K))$.

For an arbitrary cutoff $N$, take $K=\lfloor N^\alpha\rfloor$. The final incomplete block has $O_\alpha(K^{\beta-1}+1)=O_\alpha(N^{1-\alpha}+1)$ integers. Combining (1)--(3), and absorbing that incomplete block on both sides, gives
\[
 \#\{n\le N:\gcd(n,\lfloor n^\alpha\rfloor)=1\}
 =\frac6{\pi^2}N+
 O_\alpha\bigl((N^\alpha+N^{1-\alpha})\log(2N)\bigr).
\tag{4}
\]
Both error exponents are less than one. This proves the full density statement in the sublinear range, with an explicit elementary error estimate.

\section*{The precise external input for larger powers}
Bergelson--Richter's Theorem 1 states: if $f$ belongs to a Hardy field and
\[
 \log t\,\log_4t=o(f(t)),\qquad
 t^{j-1}=o(f(t)),\qquad f(t)=o(t^j)
\tag{5}
\]
for some positive integer $j$, then the natural density of $\{n:\gcd(n,\lfloor f(n)\rfloor)=1\}$ is $6/\pi^2$. Here $\log_4$ means four iterated logarithms, as in that paper. A Hardy field is a field of germs at infinity closed under differentiation.

For $f(t)=t^\alpha$ with positive noninteger $\alpha$, choose $j=\lceil\alpha\rceil$. Then $j-1<\alpha<j$, and each of the three ratios in (5) tends to zero. Powers $t^\alpha=\exp(\alpha\log t)$ belong to the logarithmico-exponential Hardy field. This verifies every hypothesis, including the noninteger restriction, and proves the result for $\alpha>1$ as well.

The restriction is necessary: if $\alpha$ is a positive integer then $\lfloor n^\alpha\rfloor=n^\alpha$, so the gcd equals $n$ and the density of coprime values is zero. Merely knowing fractional-part equidistribution for each fixed modulus would not by itself justify removing arbitrarily large common prime divisors; that extra analytic work is contained in the declared theorem.

\section*{Assessment}
The result is complete using the stated Hardy-field theorem for superlinear exponents. Equation (4) independently proves the entire sublinear range, and the exceptional integer exponents are checked directly. No new density theorem is claimed.

\textbf{Completion rate estimate: 100\%.}

\end{document}
